\section{C2 pitchfork with three arbitrary event weights}
\label{sec:c2-three-large}

\begin{theorem}[All scales of the three retained events]
\label{thm:c2-three-large}
Let $K_c,K_t$ be independent real closed extensions of $\mathbb R$.
Let a physical canonical C2 source over $K_c$ have $\lambda=\beta=0$,
with $d,C,\alpha$ positive bounded units and
$n,m,A,D\geq0$ arbitrary. Suppose an actual finite noncommuting
stationary zero has $u_c=k>0$ infinitesimal and $v_c=a>0$ a bounded
unit. Then the entire original parent is tensor SOS at every
independent projective measurement pair and every real or complex
projective state. No ratios among $n,A,D$ are restricted.
\end{theorem}

The bounded-$A$ result is a special case. The mechanism is to keep
all three possibly large event amplitudes as state coordinates.
The exact parent and its amplitudes are (TD1); in particular
$J=uW+uvX-Y-vZ$.

\paragraph{Source entry with three large weights.}
The initial coercivity estimate in the preceding proof uses only
$d,C,\alpha$, not $A$. Thus after $W_c=1$,
\[
 X_c,Y_c,Z_c=O(k),\qquad m=O(k^2).
\]
Now use $F_Z=0$ first. The $A$ and $n$ packets do not enter this
state equation. Every term other than $-8DvJ$ is $O(k)$, because
$v_c=a$ is a unit and the just displayed bounds hold. Therefore
$DJ_c=O(k)$. The $Y$ equation then gives $AY_c=O(k)$, and the $X$
equation gives $(d+2n)X_c=O(k^2)$. Write
\[
 N=d+2n,\quad T_A=d+2A,\quad L=d+2D,\qquad
 x=X_c/k^2,\quad y=Y_c/k,\quad j=J_c/k.
\]
Consequently
\[
 x,y,j,Nx,Ay,Dj\ \hbox{are bounded},\qquad N^{-1}\ \hbox{bounded}.
                                                               \tag{TA1}
\]
These estimates do not yet invert $A$ or $D$.

Use the exact stationary equations (TD3) and divide both by
$T_A L$. Now $t_2/T_A,t_6/T_A,t_3/L,t_4$ are positive bounded
units. In the first equation the coefficient multiplying $U^2$ is
$(A/T_A)$ times a positive bounded unit: its remaining bracket is
\[
 4t_4\left\{\theta\frac{t_3}{L}
            +CV\frac{1+2D/N}{L}\right\},
\]
whose lower bound follows from (TD4). The coefficient multiplying
$m$ is also a positive bounded unit:
\[
 4V\left\{\alpha\frac{t_8}{T_A L}
              +\frac{D}{L}\frac{t_6}{T_A}V\right\}.             \tag{TA2}
\]
Indeed, if $D/L$ is bounded below, the second summand is a unit.
Otherwise $D$ is infinitesimal and $L$ is a unit, while
$t_8/T_A=1+2(C+D)/T_A$ is a unit. All displayed factors have
bounded upper bounds in either case. Hence
\[
 \frac{m}{U^2}=\frac{A}{T_A}\cdot(\hbox{positive bounded unit}). \tag{TA3}
\]
The second stationary equation becomes
\[
 \frac{A}{T_A}\cdot(\hbox{positive unit})
       +O\left(\frac{A}{T_A}U\right)
 =\frac{D}{L}\cdot(\hbox{positive unit})
       +O\left(\frac{A}{T_A}U\right).                          \tag{TA4}
\]
Thus $A/T_A$ and $D/L$ are comparable. If they were infinitesimal,
$A,D,m$ would all be infinitesimal, contradicting exactly the
physicality test printed after (TD5). Therefore
\[
 A/T_A,\ D/L,\ m/k^4\quad\hbox{are positive bounded units}.       \tag{TA5}
\]
In particular $A,D$ have bounded reciprocals, but may be arbitrarily
large. Set $\gamma=4ma^2/k^4$, also a positive bounded unit.

\paragraph{The three-packet graph.}
On the local testing domain $\bar v=\bar a\ne0$, put
$b=1-j-y$ and use
\[
 X_g=xu^2W,\qquad Y_g=yuW,\qquad J_g=juW,\qquad
 Z_g=u\left(\frac b v+xu^2\right)W.                            \tag{TA6}
\]
This graph matches the whole actual contact: the definition of $J_c$
gives $Z_c/k=b/a+xk^2$. Its three hard normals are
\[
 E_X=\sqrt N(X-xu^2W),\quad
 E_Y=\sqrt A(Y-yuW),\quad E_J=\sqrt D(J-juW).                    \tag{TA7}
\]
The exact inverse is
\[
 X=xu^2W+E_X/\sqrt N,\quad Y=yuW+E_Y/\sqrt A,
\quad Z=\frac{uW+uvX-Y-juW-E_J/\sqrt D}{v}.                    \tag{TA8}
\]
Every denominator is either pure coefficient or the pure testing
unit $v$. All graph coefficients and all scaled normal coefficients
are bounded. For the three large packets this follows explicitly
from
\[
 n x^2,\ nx,\ n/N,\quad A y^2,\ \sqrt A\,y,\quad
 D j^2,\ \sqrt D\,j\quad\hbox{bounded},                       \tag{TA9}
\]
using (TA1),(TA5). Their quadratic normal contributions are exactly
$4(n/N)D_uD_v E_X^2$, $4D_uD_vE_Y^2$, and $4E_J^2$.
Every other original source coefficient is bounded, including $m$.

For clarity, the finite Taylor division used below works in the
bounded Laurent polynomial ring in $v$: for positive integer $r$,
\[
 v^{-r}-a^{-r}
 =-(v-a)\frac{\sum_{i=0}^{r-1}v^i a^{r-1-i}}{a^rv^r}.
\]
Thus a Laurent polynomial vanishing at $(U,v)=(K,a)$ belongs to
$(U-K,v-a)$ with bounded Laurent coquotients; repeating the same
identity gives the second-order division. There is no localization
at a difference of the two fields.

Set $K=k^2$, $R=u^2-(K/a)v$ and $T=u(v-a)$. The entire graph
scalar is again $4mv^2+u^2H_0(u^2,v)$, now with bounded Laurent
coefficients in $v$. Its three scalar contact jets give (TD10),
and hence an exact bounded Gram on $(R,T)$ with $R$ residue pivot
$\bar\gamma$. The full first normals before their pure source
scale divisions are
\[
\begin{split}
 (F_X+uF_Z)|_g&=-8mv+u^2L_X(u^2,v),\\
 (F_Y-F_Z/v)|_g&=uL_Y(u^2,v),\\
 (-F_Z/v)|_g&=uL_J(u^2,v).                                     \end{split} \tag{TA10}
\]
Divide the rows by $\sqrt N,\sqrt A,\sqrt D$, respectively.
The scaled coefficients are bounded by (TA9). At contact they
vanish. Laurent Taylor division gives exactly the root membership
and infinitesimal $R$ couplings proved in (TD11)--(TD12).
The $X$ row has infinitesimal coefficients on both $R,T$ because
$m/K=O(K)$; the other rows have infinitesimal coefficients on $R$.
The exact quadratic state remainder is bounded by (TA8)--(TA9),
so the entire original parent has a bounded five-root Gram on
$(E_X,E_Y,E_J,R,T)$.

\paragraph{Uniform strictness without a largest-event choice.}
The $X$ and $R$ residue pivots are $2(1+\bar a^2)$ and
$\bar\gamma$. The remaining hard state form is still (TD13).
At $u=0$, (TA8) gives
\[
 Y=E_Y/\sqrt A,\qquad
 Z=-\frac{E_Y}{v\sqrt A}-\frac{E_J}{v\sqrt D}.
\]
In these coordinates (TD13) is the sum of the two explicit pivots
\[
 4(1+v^2)E_Y^2+4E_J^2
\]
and the nonnegative bounded form
$(4C+2d)(1+v^2)Y^2+4\alpha(vY-Z)^2$.
Its residue is therefore positive definite for every combination of
bounded or unbounded $A,D$.

The boundary Schur function is the exact expression
\[
 f(v)=\frac{8(1+v^2)}{\det B(v)/(T_A L)}
 \left(\frac{Q_{10}}{T_A L}+
       \frac{Q_0}{T_A L}v^2+\frac{Q_{12}}{T_A L}v^4\right),     \tag{TA11}
\]
with $Q_{10},Q_{12}$ as in (TD14). For $v$ a bounded unit all its
coefficients and denominator are bounded, and
\[
 \frac{Q_{10}}{T_A L}=4\alpha\frac A{T_A}\frac{t_3}{L},\qquad
 \frac{Q_{12}}{T_A L}=4C\frac D L\frac{t_2-2m}{T_A}
\]
are positive units. The determinant has the positive term
$16AD(1+v^2)v^2$, so its quotient by $T_A L$ has a positive lower
bound. Its other terms are bounded after that division.
The exact middle coefficient is polynomial of degree at most one
in each of $A,D$, so $Q_0/(T_A L)$ is bounded too.

The contact normal equations and scalar jets give
$\bar f(\bar a)=\bar f'(\bar a)=0$ as before. Therefore the real
quartic bracket equals
$\overline{Q_{12}/(T_A L)}(v^2-\bar a^2)^2$.
Its strict second derivative gives the positive soft Schur pivot on
$T$. The whole five-root residue Gram is strict. This proves the
local all-state certificate; the Bob sign identity gives its negative
sign copy. Homogenization in $W$ includes $W=0$.

\paragraph{Outer tests when both $A$ and $D$ are large.}
For $u$ infinitesimal and $v$ a bounded unit with
$\bar v\ne\pm\bar a$, retain the entire $m,n$ packets and use the
four roots
\[
 X,\quad K_A=\sqrt A Y,\quad K_D=\sqrt D J,\quad uW.
\]
Substitute $Y=K_A/\sqrt A$ and
$Z=(uW+uvX-K_A/\sqrt A-K_D/\sqrt D)/v$.
The exact remainder is a bounded quadratic Gram. Its hard residue
has the positive $X$ pivot $2\bar d(1+\bar v^2)$ and the two
explicit pivots just displayed; its soft Schur complement is
$\bar f(\bar v)>0$ by (TA11).

At $u,v$ both infinitesimal, and at $u,1/v$ both infinitesimal,
one need not divide by $v$. Cap $C,A,D,\alpha$ at the positive
unit $\mu=\min(1,C/d,A/d,D/d,\alpha/d)$ after division by $d$,
and retain all excesses and the full $m,n$ packets. The bounded
selected parent has an exact Gram on $(X,Y,Z,uW)$ in the finite
chart and on the same four roots in the homogeneous reciprocal
chart. At $(u,v)=(0,0)$ its soft Schur complement is
$8Q_{10,\mu}/\det B_\mu(0)>0$; at $(u,1/v)=(0,0)$ it is
$8Q_{12,\mu}/\det B_{\mu,\infty}>0$. The hard blocks are positive
by (TD13) and its homogeneous limit. Hence both entire endpoint
neighborhoods have strict-residue Grams, including literal zeros
and infinities. These expressions use only the bounded selected
parent; no residue of an unbounded $A$ or $D$ is taken.

All other commuting tests use the same cap and the already printed
open-axis and corner forms in the preceding all-tests proof. Their
hypotheses require only the four lower bounds supplied by (TA5).
Noncommuting residue tests use the retained arbitrary-physical-source
theorem. Restore the positive pure factors and split complex states
into real and imaginary parts. This completes every independent
test for the stated source domain.

\paragraph{Coverage gained and remaining restriction.}
The theorem removes every relative-scale restriction among $n,A,D$
from the finite pitchfork. Its reflected form covers the corresponding
ordinary--infinite actual contacts. It still assumes $C/d$ and
$\alpha/d$ are bounded units and this actual-angle residue type.
Unbounded $C/d$ or $\alpha/d$, other actual-angle types, and the
exhaustive C2 join remain outside this theorem.
